\documentclass[10pt,twoside]{article}

\usepackage[utf8]{inputenc}
\usepackage[T1]{fontenc}
\usepackage[letterpaper,textwidth=12.2cm,textheight=19.3cm,top=3.0cm,asymmetric]{geometry}
\usepackage{amsmath,amssymb}
\usepackage{booktabs}
\usepackage{array}
\usepackage{graphicx}
\graphicspath{{figures/}{paper/figures/}}
\usepackage{xcolor}
\usepackage{tikz}
\usepackage{url}
\usepackage{enumitem}
\usepackage{float}
\usepackage[font=small,labelfont=bf,labelsep=period]{caption}
\usepackage{fancyhdr}
\usepackage[colorlinks=true,citecolor=blue,urlcolor=blue,linkcolor=blue]{hyperref}
\usetikzlibrary{arrows.meta,positioning,calc,fit,backgrounds}

\renewcommand{\figurename}{Fig.}
\setlength{\parindent}{1em}
\setlength{\parskip}{0pt}
\setlist[itemize]{label=\textendash,leftmargin=1.2em,itemsep=0.1em,topsep=0.25em}

\definecolor{paperBlue}{RGB}{37,99,235}
\definecolor{paperCyan}{RGB}{6,182,212}
\definecolor{paperGreen}{RGB}{22,163,74}
\definecolor{paperOrange}{RGB}{234,88,12}
\definecolor{paperInk}{RGB}{15,23,42}
\definecolor{paperGray}{RGB}{100,116,139}
\definecolor{paperLight}{RGB}{248,250,252}
\definecolor{paperBorder}{RGB}{203,213,225}

\newcommand{\method}{MixiSpeech}
\newcommand{\mat}{MAT}
\newcommand{\rtf}{RTF}
\newcommand{\xfinal}{x_{\mathrm{final}}}
\newcommand{\xhat}{\hat{x}_{\mathrm{final}}}
\newcommand{\best}[1]{\textbf{#1}}
\newcommand{\pending}{\textit{TBD}}

\pagestyle{fancy}
\fancyhf{}
\fancyhead[LE]{\small\thepage\hspace{2.0em}MixiSpeech Team}
\fancyhead[RO]{\small\method: One-Step TTS Distillation\hspace{2.0em}\thepage}
\renewcommand{\headrulewidth}{0pt}

\begin{document}
\thispagestyle{empty}

\begin{center}
\vspace*{0.45cm}
{\LARGE\bfseries
\method: Progressive One-Step Distillation\\
for Vietnamese--English Text-to-Speech\par}
\vspace{0.85cm}
{\normalsize
MixiSpeech Team\textsuperscript{1}\\}
\vspace{0.35cm}
{\small
\textsuperscript{1}MixiSpeech Research, Hanoi, Vietnam\\
\texttt{mixispeech@mixiuniversal.com}}
\end{center}

\vspace{0.65cm}

\begin{center}
\begin{minipage}{0.88\textwidth}
\small
\noindent\textbf{Abstract.}
Diffusion and flow-matching text-to-speech (TTS) systems produce high-quality
speech, but iterative sampling makes deployment expensive for interactive
applications. We present \method, a compact one-step student distilled from a
multi-step Matcha-TTS teacher for Vietnamese--English speech synthesis. The core
contributions are: (1) a trajectory-aware one-step distillation objective that
combines final-state, velocity, and local pair-consistency supervision; (2) a
progressive pair-stride curriculum
$64 \rightarrow 32 \rightarrow 16 \rightarrow 8 \rightarrow 4 \rightarrow 2
\rightarrow 1$ that transfers a 64-step teacher trajectory into a single
student update; and (3) an empirical study showing that trajectory
distillation recovers quality lost by naive one-step Matcha sampling while
retaining one-step latency. On FLEURS Vietnamese evaluation, \method{} keeps the
same 20.94M-parameter footprint and latency class as a naive one-step Matcha
teacher, but improves UTMOS from $1.575$ to $2.264$, WER from $16.22\%$ to
$12.84\%$, and CER from $8.15\%$ to $5.87\%$. Across 200 Vietnamese FLEURS
prompts and 200 English VCTK prompts, the model uses one acoustic step and is
substantially faster than 32/64-step teacher sampling. For reproducibility, we
describe the Vietnamese--English data-processing pipeline and provide public
evaluation scripts; audio samples and a live demo are available at
\url{https://mixispeech.solanai.us/}.

\vspace{0.55em}
\noindent\textbf{Keywords:}
Text-to-speech $\cdot$ Flow matching $\cdot$ Progressive distillation $\cdot$
Vietnamese speech synthesis $\cdot$ Low-latency inference.
\end{minipage}
\end{center}

\section{Introduction}

\begingroup
\setlength{\parindent}{0pt}
\setlength{\parskip}{0.55em}

Modern neural TTS systems often trade inference latency for quality. Diffusion
and flow-based acoustic models~\cite{ddpm,flow_matching} refine a latent or mel representation through
many sampling steps, which is acceptable for offline generation but costly for
interactive applications. A tempting deployment shortcut is to reduce the number
of solver steps directly. {\color{red}In our setting, however, the one-step Matcha teacher
has the same 20.94M-parameter footprint and similar runtime as the student, but
loses perceptual quality and intelligibility on Vietnamese evaluation.}

\method{} addresses this gap with one-step distillation. Instead of supervising
the student only on the final teacher sample, we distill dense teacher
trajectories and combine final-state consistency, velocity matching, and local
pair-consistency supervision. The resulting student performs a single acoustic
update at inference time but learns from the geometry of the teacher's
multi-step path.

This paper makes five contributions. First, we show that naive one-step teacher
sampling is not an adequate deployment baseline for Vietnamese technical and
general speech. Second, we propose a progressive trajectory-consistency recipe
that improves one-step quality without increasing inference steps. Third, we
document the implemented composite-corpus pipeline and Matcha fine-tuning recipe
used to produce the teacher and student checkpoints. Fourth, we evaluate
\method{} against \mat-1, \mat-4, \mat-32, \mat-64, VieNeu-TTS, and F5-TTS on
public Vietnamese and English evaluation data. Finally, we define a release
boundary for the composite corpus, public prompt lists, objective-evaluation
scripts, demo audio, and a blinded human listening protocol so the work can be
independently audited without redistributing third-party audio outside its
source licenses.

This work is positioned relative to non-autoregressive and flow-based TTS
systems such as FastSpeech~2, Glow-TTS, VITS, Matcha-TTS, Voicebox, E2 TTS, and
F5-TTS~\cite{fastspeech2,glowtts,vits,matcha,voicebox,e2tts,f5tts}; iterative
diffusion/flow acoustic models and their fast samplers
\cite{ddpm,ddim,gradtts,difftts,prodiff}; and consistency/progressive
distillation methods that motivate one-step generation
\cite{progressive_distillation,consistency}. Unlike large zero-shot systems
such as VALL-E~\cite{valle}, our goal is a small, auditable one-step acoustic
student for the Vietnamese--English deployment setting.

\endgroup

\section{Method}

\subsection{Teacher and Student}

The teacher is a Matcha-TTS acoustic model fine-tuned for the MixiSpeech
setting. At inference, the teacher can be sampled with different ODE/flow steps.
We evaluate \mat-1, \mat-4, \mat-32, and \mat-64 to characterize the
quality--latency trade-off of the original teacher.

The student uses the same Matcha-family architecture footprint as the teacher,
but is trained to predict the final teacher mel state from an intermediate
teacher trajectory state. Let $x_t$ be a trajectory state at time $t$,
$\xfinal$ be the final teacher state, and $c$ denote text, duration, and speaker
conditioning. The student predicts a velocity field $v_\theta$ and forms a
one-step final estimate:
\begin{equation}
  \xhat = x_t + (1 - t)\, v_\theta(x_t, t, c).
  \label{eq:jump}
\end{equation}

\subsection{Progressive Pair-Consistency Distillation}

The objective combines three terms:
\begin{align}
\mathcal{L}_{\mathrm{final}} &= \|\xhat - \xfinal\|_1, \\
\mathcal{L}_{\mathrm{vel}} &= \left\|v_\theta(x_t,t,c) -
  \frac{\xfinal - x_t}{1-t}\right\|_2^2, \\
\mathcal{L}_{\mathrm{pair}} &= \|\Phi_\theta^{(r)}(x_i,t_i,c) - x_{i+r}\|_1, \\
\Phi_\theta^{(r)}(x_i,t_i,c) &= x_i + (t_{i+r}-t_i)\,
  v_\theta(x_i,t_i,c),
\end{align}
where $(x_i,x_{i+r})$ is a teacher-state pair selected at the current curriculum
stride $r$, and $\Phi_\theta^{(r)}$ denotes the local student update toward the
paired target state. The total loss is
\begin{equation}
  \mathcal{L} =
  \lambda_f\mathcal{L}_{\mathrm{final}} +
  \lambda_v\mathcal{L}_{\mathrm{vel}} +
  \lambda_p\mathcal{L}_{\mathrm{pair}}.
\end{equation}

Dense 64-step teacher trajectories are distilled with progressively shorter pair
strides:
\begin{equation}
  64 \rightarrow 32 \rightarrow 16 \rightarrow 8 \rightarrow 4 \rightarrow 2
  \rightarrow 1.
\end{equation}
The schedule is inspired by progressive distillation and consistency-style
learning, but the implemented recipe is a local hybrid: it retains direct
final-state supervision while exposing the student to local trajectory pairs.
{\color{red}The best checkpoint in the current experiments is selected from the progressive
run around the stride-2 stage and remains one-step at inference.}

\begin{figure}[t]
\centering
\includegraphics[width=\textwidth]{arch.png}
\caption{\method{} architecture and training flow. A frozen 64-step
Matcha-TTS teacher generates dense acoustic trajectories offline. The one-step
student is trained with final-state, velocity-matching, and local
pair-consistency losses under a progressive stride curriculum. In the pair loss,
$\Phi_\theta$ denotes a local stride update with $\Delta t=t_{i+r}-t_i$, while
the final-state jump uses Eq.~\eqref{eq:jump}. At inference, only the trained
student and frozen HiFi-GAN vocoder~\cite{hifigan} are used, so acoustic
generation takes one student step.}
\label{fig:architecture}
\end{figure}

\section{Implemented Data and Training Pipeline}

The paper prototype follows the same pipeline used by the current local
checkpoints. We describe it explicitly because it is easy to overstate the
method if the teacher is treated as an abstract pretrained model. The current
system first constructs a Matcha-compatible bilingual corpus, then fine-tunes a
conditional-flow teacher, then extracts teacher trajectories for the student.

\subsection{Audio Audit and Bilingual Manifests}

All audio is standardized to 24 kHz and converted into Matcha filelists of the
form \texttt{wav|speaker\_id|phonemes}. The dataset used here is not a single
public corpus and is not introduced under an internal run name. It is a
composite training filelist built from named sources. The starting
Vietnamese/code-switch recipe in the previous draft contained audited VieNeu
speech, PhoAudio Vietnamese audiobook speech, synthetic Vietnamese--English
code-switch utterances, and a small Vietnamese math/chemistry subset. That
earlier recipe contains 95,757 training utterances (169.79 h) over 209 speakers,
with manifest components VieNeu 70,258; PhoAudio 13,073; code-switch 11,749; and
technical 677. The current teacher corpus uses this data design as its baseline
and adds the additional data used in this work: LibriTTS-R English speech and a
small VietMix code-switch augmentation. For reproducibility,
Table~\ref{tab:composite_data}
reports both the earlier recipe summary and the active manifests used by the
current experiments.

The upstream VieNeu audit keeps 73,524 of 73,882 rows over 193 speakers. The
rejection rules cover long leading or trailing silence, low speech ratio,
invalid duration, and frontend incompatibility; rejected items are written to
audit summaries instead of being silently mixed into training. LibriTTS-R
\texttt{train-clean-100} and \texttt{train-clean-360} contribute the English
side~\cite{libritts,librittsr} and are capped at about 100 training hours to
avoid overwhelming the Vietnamese target domain. The text frontends are
language-specific before they enter Matcha: Vietnamese uses the VieNeu IPA
frontend, and English uses
\texttt{english\_cleaners2} followed by the same IPA simplifier. The resulting
speaker map has \texttt{n\_spks}=1348 and the mel statistics are fixed at mean
$-5.5434$ and standard deviation $2.5716$.

\begin{table}[H]
\centering
\caption{Composite Matcha-compatible filelists used by the teacher
fine-tuning pipeline. The first row spells out the earlier Vietnamese/code-switch
recipe from the previous draft; the remaining rows report the current bilingual
teacher manifests used in this paper. Hours are computed from the processed
manifests and code-switch filelists.}
\label{tab:composite_data}
\scriptsize
\setlength{\tabcolsep}{3.8pt}
\resizebox{\textwidth}{!}{%
\begin{tabular}{@{}llrrrl@{}}
\toprule
Split & Component & Rows & Hours & Speakers & Notes \\
\midrule
Reference & Previous Vietnamese/code-switch recipe & 95,757 & 169.79 & 209 & VieNeu, PhoAudio, code-switch, technical \\
\midrule
Train & Current Vietnamese side & 84,008 & 152.42 & 209 & VieNeu 70,258; PhoAudio 13,073; technical 677 \\
Train & LibriTTS-R train-clean-100/360 & 61,931 & 100.00 & 1,128 & English cleaners + IPA simplifier \\
Train & VietMix code-switch augment & 561 & 0.93 & 1 & VieNeu v3 Turbo, speaker id 74 \\
\midrule
Train & Final composite teacher filelist & 146,500 & 253.35 & 1,348 & Vietnamese + English + VietMix \\
\midrule
Validation & Base bilingual & 1,516 & 2.62 & 587 & Vietnamese + English validation \\
Validation & VietMix code-switch augment & 30 & 0.05 & 1 & same frontend and speaker id \\
Validation & Final composite teacher filelist & 1,546 & 2.67 & 587+ & Vietnamese + English + VietMix \\
\bottomrule
\end{tabular}}
\end{table}

\subsection{Dataset Positioning and Release Boundary}

For an international submission, we do not position the composite training
filelist as a new benchmark dataset. It is a training recipe assembled from
named Vietnamese, English, technical, and code-switch sources, and some source
audio may be redistributable only through its original provider. The reproducible
part of the contribution is therefore the processing boundary: audit rules,
frontend normalization, filelist construction, prompt lists, evaluation scripts,
and demo audio generated by the evaluated systems.

\begin{table}[H]
\centering
\caption{Dataset and artifact release plan. The submission separates training
data used to fit the teacher from public evaluation prompts and auditable
artifacts. This avoids claiming a proprietary or partly synthetic training mix as
a public benchmark while still making the evaluation path reproducible.}
\label{tab:release_boundary}
\footnotesize
\setlength{\tabcolsep}{3.0pt}
\begin{tabular}{@{}p{0.22\textwidth}p{0.20\textwidth}p{0.28\textwidth}p{0.21\textwidth}@{}}
\toprule
Artifact & Role in paper & Public/review release & License boundary \\
\midrule
Vietnamese manifests & VieNeu, PhoAudio, and technical teacher training & derived manifests, audit stats, cleaners & audio follows source-provider terms \\
LibriTTS-R subset & English teacher training & subset recipe and filelist builder & public corpus citation and license \\
VietMix code-switch text & augmentation and stress prompts & accepted/rejected text IDs, generator recipe & text source and TTS-provider terms \\
FLEURS \texttt{vi\_vn}, VCTK prompts & objective evaluation only & exact 200+200 prompt IDs and scripts & public benchmark data; no train overlap \\
Demo and human-study audio & qualitative inspection & hosted demo, anonymized listening scores & generated outputs, not raw third-party audio \\
\bottomrule
\end{tabular}
\end{table}

This framing keeps the paper's empirical claim focused: \method{} is evaluated
on public held-out prompts and human review audio, while the training corpus is
described as an engineering recipe rather than as a new dataset contribution.
For a camera-ready version, the prompt IDs, system outputs used in human tests,
and anonymized listener scores should be archived with the demo page.

\subsection{Code-Switch Augmentation}

The current run does not replace the corpus with fully LLM-generated text.
Instead, LLM-generated sentences are reserved for hard-case prompting and
frontend stress tests. The added training augmentation used here comes from the
public
\texttt{razent/vietmix} split, filtered to short Vietnamese--English
code-switch utterances. From 1,002 candidate texts, 591 are accepted and
synthesized with VieNeu-TTS v3 Turbo using the \textit{Truc Ly} voice. The
synthetic audio is constrained to 0.35--18.0 s, resampled to 24 kHz, phonemized
with the same Vietnamese/English mixed frontend, and assigned to an existing
speaker slot, \texttt{capybara1812\_1056} (\texttt{speaker\_id}=74). The split
adds 561 utterances to teacher training and 30 to validation.

This augmentation is intentionally small compared with the bilingual base. Its
role is not to dominate the acoustic distribution, but to expose the Matcha
teacher to the code-switch and technical-token patterns that appeared in manual
testing, such as Latin abbreviations, English product words inside Vietnamese
sentences, and formula-like text after frontend normalization.

\subsection{Matcha Teacher Fine-Tuning}

The teacher is trained with the composite Vietnamese--English Matcha
configuration with added code-switch rows, summarized in
Table~\ref{tab:composite_data}. The acoustic
representation uses 80 mel bins, FFT size 1024, hop length 256, window length
1024, and frequency range 0--8 kHz. Training uses the IPA simplifier cleaner,
blank-token insertion, batch size 8, eight workers, and no precomputed duration
loading. The model keeps the Matcha vocabulary size at 187 and output size at
172.

Fine-tuning is initialized strictly from the previous bilingual teacher
checkpoint. The current composite-corpus experiment validates every 2,000
optimizer steps, logs every 50 steps, and saves the top three checkpoints every
2,000 steps with a 50,000-step training budget. The active teacher used for
trajectory extraction and evaluation in this paper is the last checkpoint from
the run started on 15 July 2026 at 02:41:25.

\subsection{Trajectory Dump and Student Run}

After teacher fine-tuning, the current student run dumps dense teacher
trajectories only for the added VietMix code-switch subset, not for the entire
composite teacher filelist. The dump uses the active composite-corpus teacher
with 64 flow steps, all state indices 0--63, temperature 0.667, and length scale
1.0. It produces 560 usable training trajectories and 30 validation
trajectories; each item stores the source waveform path, speaker id, text,
mel-frame metadata, teacher step count, state times, and serialized trajectory
tensors.

Student training then follows the progressive schedule
$64,32,16,8,4,2,1$. Each stage trains for 2,000 updates with batch size 4,
mixed precision on CUDA, learning rate $2\times 10^{-5}$, validation every 500
steps, and checkpointing every 1,000 steps. To keep conditioning stable and to
reduce the number of moving parts, the text encoder and speaker embedding are
frozen; only the Matcha decoder is trainable. The implemented loss weights are
$\lambda_f=1.0$, $\lambda_v=0.1$, and $\lambda_p=0.25$. The final run writes
\texttt{student\_last.ckpt}, while the best checkpoint selected by validation is
from the stride-2 stage in the same progressive run.

\section{Experimental Setup}

\subsection{Systems}

We compare \method{} against four teacher sampling regimes, \mat-1, \mat-4,
\mat-32, and \mat-64. The Matcha teacher and the \method{} student share the
same 20.94M-parameter architecture footprint. We also include VieNeu-TTS v3
Turbo, VieNeu-TTS v2, and F5-TTS as larger pretrained reference systems. These
references are useful for deployment context, but they are not parameter-matched
to \method{}.

\subsection{Evaluation Data and Metrics}

The public objective evaluation uses 200 Vietnamese prompts from FLEURS
\texttt{vi\_vn} test~\cite{fleurs} and 200 English prompts from VCTK
v0.92~\cite{vctk}. With eight
systems, the benchmark synthesizes 3,200 audio files. We report peak CPU RAM,
peak GPU RAM, average inference seconds, average real-time factor, UTMOS, WER,
and CER. Vietnamese ASR uses PhoWhisper-tiny~\cite{phowhisper}; English ASR
uses Whisper-tiny~\cite{whisper}. UTMOS uses the SpeechMOS UTMOS22
model~\cite{utmos}. WER and CER are diagnostic automatic metrics and do not
replace human MOS or ABX tests.

\subsection{Human Listening Study Protocol}

To make the result suitable for a conference submission, the automatic
evaluation is paired with a blinded human review round. The review set should
contain 60 prompts: 30 Vietnamese FLEURS prompts, 15 Vietnamese technical
normalization prompts covering numbers, chemistry, and math, and 15
Vietnamese--English code-switch prompts from the VietMix-derived stress set.
System names are hidden and audio is randomized per listener. We recommend at
least 20 Vietnamese listeners, with each utterance rated by at least five
listeners after balancing for fatigue. The form asks for 5-point naturalness MOS,
5-point intelligibility MOS, and AB preference against the latency-matched
\mat-1 baseline. The final paper should report mean $\pm$ 95\% confidence
intervals and a paired preference test.

\begin{table}[H]
\centering
\caption{Blinded human listening-study design for the review round. The score
columns in Table~\ref{tab:human_eval_template} are intentionally left pending in
this draft and should be filled only after collecting listener ratings.}
\label{tab:human_eval_protocol}
\scriptsize
\setlength{\tabcolsep}{4.0pt}
\resizebox{\textwidth}{!}{%
\begin{tabular}{@{}llrrl@{}}
\toprule
Prompt group & Purpose & Prompts & Systems & Human metrics \\
\midrule
FLEURS \texttt{vi\_vn} & public Vietnamese naturalness/intelligibility & 30 & 4--5 & MOS, intelligibility, AB preference \\
Technical Vietnamese & frontend robustness for math/chemistry text & 15 & 4--5 & MOS, error flag, AB preference \\
VietMix code-switch & Vietnamese--English mixed utterances & 15 & 4--5 & MOS, code-switch correctness, AB preference \\
\bottomrule
\end{tabular}}
\end{table}

\section{Results}

\begin{table}[H]
\centering
\caption{Public objective evaluation. Each model is evaluated on 200
Vietnamese prompts and 200 English prompts. Arrows in column headers indicate
the preferred direction. Bold marks the best value within each evaluation split;
ties are bolded.}
\label{tab:public_eval}
\scriptsize
\setlength{\tabcolsep}{2.7pt}
\resizebox{\textwidth}{!}{%
\begin{tabular}{@{}llrrrrrrrr@{}}
\toprule
Model & Benchmark & Params $\downarrow$ & CPU $\downarrow$ & GPU $\downarrow$ & Infer $\downarrow$ & \rtf{} $\downarrow$ & UTMOS $\uparrow$ & WER $\downarrow$ & CER $\downarrow$ \\
& & (M) & GiB & GiB & s & & & \% & \% \\
\midrule
\textbf{\method-1} & FLEURS vi\_vn & \best{20.94} & 2.41 & \best{1.38} & \best{0.040} & \best{0.005} & 2.264 & 12.84 & 5.87 \\
\mat-1 teacher & FLEURS vi\_vn & \best{20.94} & 2.41 & 1.49 & 0.041 & \best{0.005} & 1.575 & 16.22 & 8.15 \\
\mat-4 teacher & FLEURS vi\_vn & \best{20.94} & 2.41 & 1.49 & 0.056 & 0.006 & 2.143 & 17.62 & 9.08 \\
\mat-32 teacher & FLEURS vi\_vn & \best{20.94} & 2.42 & 1.49 & 0.191 & 0.022 & 2.188 & 18.82 & 9.49 \\
\mat-64 teacher & FLEURS vi\_vn & \best{20.94} & 2.47 & 1.49 & 0.343 & 0.040 & 2.194 & 18.86 & 9.52 \\
VieNeu-TTS v3 Turbo & FLEURS vi\_vn & 152.88 & \best{2.39} & 2.01 & 1.310 & 0.189 & 2.530 & \best{12.24} & \best{5.43} \\
VieNeu-TTS v2 & FLEURS vi\_vn & 293.67 & 3.33 & 1.90 & 2.147 & 0.330 & 2.767 & 13.07 & 5.73 \\
F5-TTS & FLEURS vi\_vn & 337.10 & 2.89 & 2.63 & 0.190 & 0.024 & \best{3.040} & 14.35 & 6.74 \\
\midrule
\textbf{\method-1} & VCTK v0.92 & \best{20.94} & 2.41 & \best{1.38} & \best{0.018} & \best{0.005} & 3.560 & 8.04 & 3.06 \\
\mat-1 teacher & VCTK v0.92 & \best{20.94} & 2.41 & 1.49 & 0.019 & \best{0.005} & 3.491 & 2.72 & 0.76 \\
\mat-4 teacher & VCTK v0.92 & \best{20.94} & 2.41 & 1.49 & 0.029 & 0.009 & 4.160 & 3.08 & 0.91 \\
\mat-32 teacher & VCTK v0.92 & \best{20.94} & 2.42 & 1.49 & 0.121 & 0.035 & 4.179 & 3.52 & 1.07 \\
\mat-64 teacher & VCTK v0.92 & \best{20.94} & 2.47 & 1.49 & 0.228 & 0.066 & 4.177 & 3.60 & 1.13 \\
VieNeu-TTS v3 Turbo & VCTK v0.92 & 152.88 & \best{2.39} & 2.01 & 0.757 & 0.188 & 4.136 & 4.22 & 1.44 \\
VieNeu-TTS v2 & VCTK v0.92 & 293.67 & 3.33 & 1.90 & 2.581 & 0.335 & 3.615 & 25.74 & 105.60 \\
F5-TTS & VCTK v0.92 & 337.10 & 2.89 & 2.63 & 0.207 & 0.035 & \best{4.374} & \best{1.62} & \best{0.37} \\
\bottomrule
\end{tabular}}
\end{table}

\begin{table}[H]
\centering
\caption{Vietnamese target distillation ablation on 200 FLEURS \texttt{vi\_vn}
prompts. The full row is the validation-selected checkpoint used in the main
system; ablation rows are 14k-update final checkpoints trained from the same
composite-corpus teacher. Arrows indicate the preferred direction. The rightmost
column averages WER and CER to summarize target ASR error. Bold marks the best
value; ties are bolded.}
\label{tab:distill_ablation}
\scriptsize
\setlength{\tabcolsep}{3.0pt}
\resizebox{\textwidth}{!}{%
\begin{tabular}{@{}lrrrrrrrrr@{}}
\toprule
Ablation & N & Params $\downarrow$ & CPU $\downarrow$ & GPU $\downarrow$ & Infer $\downarrow$ & \rtf{} $\downarrow$ & WER $\downarrow$ & CER $\downarrow$ & Avg. ASR $\downarrow$ \\
& & (M) & GiB & GiB & s & & \% & \% & error \% \\
\midrule
Full: progressive + $\mathcal{L}_{pair}$ (ours) & 200 & \best{20.94} & \best{2.41} & \best{1.38} & \best{0.040} & \best{0.005} & \best{12.84} & \best{5.87} & \best{9.36} \\
w/o $\mathcal{L}_{pair}$ & 200 & \best{20.94} & 2.43 & \best{1.38} & \best{0.040} & \best{0.005} & 13.23 & 6.05 & 9.64 \\
w/o progressive curriculum & 200 & \best{20.94} & 2.45 & \best{1.38} & \best{0.040} & \best{0.005} & 13.22 & 6.16 & 9.69 \\
Direct $\mathcal{L}_{jump}+\mathcal{L}_{vel}$ & 200 & \best{20.94} & \best{2.41} & \best{1.38} & \best{0.040} & \best{0.005} & 13.04 & 5.94 & 9.49 \\
\bottomrule
\end{tabular}}
\end{table}

\subsection{Human Listening Review Round}

Table~\ref{tab:human_eval_template} is the score table for the listening review
round. It is deliberately not populated with synthetic numbers. The paper should
only submit human results after listener scores are collected from the protocol
in Table~\ref{tab:human_eval_protocol}; until then, the automatic metrics in
Tables~\ref{tab:public_eval}--\ref{tab:distill_ablation} should be treated as
diagnostic evidence rather than final perceptual validation.

\begin{table}[H]
\centering
\caption{Human evaluation table to be filled after blinded listening review.
Nat. MOS and Intell. MOS use a 1--5 scale with 95\% confidence intervals.
Preference is the percentage of AB comparisons preferring the listed system over
the latency-matched \mat-1 baseline. Values are pending listener collection.}
\label{tab:human_eval_template}
\scriptsize
\setlength{\tabcolsep}{3.2pt}
\resizebox{\textwidth}{!}{%
\begin{tabular}{@{}lrrrrl@{}}
\toprule
System & Steps & Prompts & Nat. MOS $\uparrow$ & Intell. MOS $\uparrow$ & AB pref. vs \mat-1 $\uparrow$ \\
\midrule
\method-1 (ours) & 1 & 60 & \pending{} & \pending{} & \pending{} \\
\mat-1 teacher & 1 & 60 & \pending{} & \pending{} & baseline \\
\mat-64 teacher & 64 & 60 & \pending{} & \pending{} & \pending{} \\
VieNeu-TTS v3 Turbo & reference & 60 & \pending{} & \pending{} & \pending{} \\
F5-TTS & reference & 60 & \pending{} & \pending{} & \pending{} \\
\bottomrule
\end{tabular}}
\end{table}

\subsection{Distillation Ablation}

Table~\ref{tab:distill_ablation} addresses whether pair consistency and the
progressive curriculum improve the target Vietnamese setting. The answer is
positive for intelligibility: the full recipe gives the best FLEURS WER, CER,
and average ASR error among the tested one-step students. Removing
$\mathcal{L}_{pair}$ increases WER from 12.84\% to 13.23\% and CER from 5.87\%
to 6.05\%. Removing the progressive curriculum while randomly sampling pair
strides gives the highest CER, 6.16\%. The direct
$\mathcal{L}_{jump}+\mathcal{L}_{vel}$ baseline remains competitive, but is
worse than the full recipe on all three Vietnamese ASR error metrics.

We do not claim that the full recipe wins every automatic diagnostic. In the
same FLEURS run, the no-progressive variant has a slightly higher UTMOS
(2.277 versus 2.264), and VCTK English diagnostics are less conclusive. On the
code-switch sets used to motivate the augmentation, however, progressive
distillation improves UTMOS over direct final-only training from 2.3959 to
2.4613 on VietMix val30 and from 2.0837 to 2.1542 on VietMix stress200, with the
same one-step inference budget. We therefore treat pair consistency and
progressive scheduling as stabilizing trajectory supervision for the target
Vietnamese/code-switch setting, not as a claim that one-step training is
impossible without them.

\subsection{Vietnamese Target Evaluation}

Table~\ref{tab:public_eval} shows that \method{} and \mat-1 occupy the same
runtime class on FLEURS Vietnamese: 0.040s versus 0.041s average inference
time, both at 20.94M parameters. However, \method{} substantially improves the
automatic quality and intelligibility metrics. UTMOS increases from 1.575 to
2.264, WER decreases from 16.22\% to 12.84\%, and CER decreases from 8.15\% to
5.87\%. This supports the main deployment claim: training a one-step student
from teacher trajectories recovers substantial Vietnamese quality that is lost
by naive one-step teacher sampling.

The multi-step teacher improves perceptual quality relative to \mat-1, but
incurs higher inference time. \mat-64 is 8.6$\times$ slower than \method{} on
FLEURS under the measured setup. \method{} therefore targets the latency point
of \mat-1 while recovering a meaningful part of the quality lost by naive step
reduction.

\subsection{English Diagnostic Evaluation}

On VCTK, the teacher-family systems remain stronger in WER/CER, and F5-TTS
reaches the highest UTMOS. This is expected because VCTK is an English corpus
and the \method{} student is optimized for the target Vietnamese--English
setting. We use VCTK as an out-of-domain English diagnostic rather than as the
primary target claim.

\section{Discussion}

\method{} is not presented as an absolute-quality replacement for much larger
pretrained TTS systems. F5-TTS and VieNeu-TTS remain strong references,
especially in English naturalness and general pretrained coverage. The
contribution is narrower and deployment-oriented: a one-step student can be
trained to be meaningfully better than naive one-step teacher sampling at the
same parameter count and latency class.

The result matters because \mat-1 is the most direct one-step baseline. If a
paper only compares a one-step student to \mat-64, the speedup is expected and
does not prove that distillation is necessary. By including \mat-1, \mat-4,
\mat-32, and \mat-64, the evaluation separates the effect of step count from the
effect of trajectory-aware student training.

\section{Limitations}

The present draft includes a human-study protocol but not yet collected human
scores. UTMOS, WER, and CER are useful for rapid iteration, but
publication-quality validation should replace Table~\ref{tab:human_eval_template}
with real MOS and preference results. WER/CER can also be biased by ASR behavior,
especially for chemical formulas, math expressions, named entities, and
code-switching. The current system is not a full reference-audio voice cloning
model; it uses fixed speaker IDs for the Matcha-family systems and
preset/reference voices for external baselines.

\section{Conclusion}

We introduced \method, a one-step Vietnamese--English TTS student distilled
from a multi-step Matcha teacher. The method combines final-state supervision,
velocity matching, and progressive pair consistency over dense teacher
trajectories. On the primary Vietnamese public evaluation, \method{} keeps the
speed and size of the one-step teacher baseline while improving UTMOS, WER, and
CER. The results indicate that progressive trajectory distillation is a
practical path toward low-latency TTS without relying solely on naive
solver-step reduction.

\paragraph{Reproducibility.}
The main shell entry points are included with the repository:
\begingroup
\footnotesize
\sloppy
\begin{itemize}[leftmargin=1.5em,itemsep=0.1em,topsep=0.2em]
\item data restore and manifest validation:
\path{experiments/run_restore_vieneu_140h_tmux.sh};
\item teacher fine-tuning:
\path{experiments/run_mix5_resolved_teacher_tmux.sh} and
\path{experiments/run_mix5_codeswitch_teacher_only_tmux.sh};
\item progressive student distillation and trajectory dumping:
\path{experiments/run_progressive_matcha_distill_tmux.sh};
\item ablation training/evaluation:
\path{experiments/run_distill_ablation_tmux.sh};
\item public objective evaluation:
\path{experiments/run_public_objective_eval.sh} with
\path{experiments/configs/objective_eval_200.env};
\item local Streamlit demo:
\path{experiments/run_streamlit_mixispeech_tts_demo.sh}.
\end{itemize}
\endgroup
\noindent
The public objective run uses \texttt{VI\_LIMIT=200}, \texttt{EN\_LIMIT=200},
FLEURS \texttt{vi\_vn}, VCTK v0.92, CUDA devices for synthesis/UTMOS, and all
systems compared in Table~\ref{tab:public_eval}.

\begin{thebibliography}{99}
\bibitem{matcha}
S. Mehta, R. Tu, J. Beskow, E. Sz\'ekely, and G. E. Henter, ``Matcha-TTS: A fast
TTS architecture with conditional flow matching,'' in \emph{Proc. ICASSP}, 2024.

\bibitem{flow_matching}
Y. Lipman, R. T. Q. Chen, H. Ben-Hamu, M. Nickel, and M. Le, ``Flow matching
for generative modeling,'' in \emph{Proc. ICLR}, 2023.

\bibitem{ddpm}
J. Ho, A. Jain, and P. Abbeel, ``Denoising diffusion probabilistic models,'' in
\emph{Proc. NeurIPS}, 2020.

\bibitem{ddim}
J. Song, C. Meng, and S. Ermon, ``Denoising diffusion implicit models,'' in
\emph{Proc. ICLR}, 2021.

\bibitem{progressive_distillation}
T. Salimans and J. Ho, ``Progressive distillation for fast sampling of diffusion
models,'' in \emph{Proc. ICLR}, 2022.

\bibitem{prodiff}
R. Huang, Z. Zhao, H. Liu, J. Liu, C. Cui, and Y. Ren, ``ProDiff: Progressive
fast diffusion model for high-quality text-to-speech,'' in \emph{Proc. ACM
Multimedia}, 2022.

\bibitem{consistency}
Y. Song, P. Dhariwal, M. Chen, and I. Sutskever, ``Consistency models,'' in
\emph{Proc. ICML}, 2023.

\bibitem{fastspeech2}
Y. Ren, C. Hu, X. Tan, T. Qin, S. Zhao, Z. Zhao, and T.-Y. Liu, ``FastSpeech 2:
Fast and high-quality end-to-end text to speech,'' in \emph{Proc. ICLR}, 2021.

\bibitem{glowtts}
J. Kim, S. Kim, J. Kong, and S. Yoon, ``Glow-TTS: A generative flow for
text-to-speech via monotonic alignment search,'' in \emph{Proc. NeurIPS}, 2020.

\bibitem{vits}
J. Kim, J. Kong, and J. Son, ``Conditional variational autoencoder with
adversarial learning for end-to-end text-to-speech,'' in \emph{Proc. ICML},
2021.

\bibitem{gradtts}
V. Popov, I. Vovk, V. Gogoryan, T. Sadekova, and M. Kudinov, ``Grad-TTS: A
diffusion probabilistic model for text-to-speech,'' in \emph{Proc. ICML}, 2021.

\bibitem{difftts}
M. Jeong, H. Kim, S. Cheon, B. J. Choi, and N. S. Kim, ``Diff-TTS: A denoising
diffusion model for text-to-speech,'' in \emph{Proc. Interspeech}, 2021.

\bibitem{voicebox}
M. Le et al., ``Voicebox: Text-guided multilingual universal speech generation
at scale,'' in \emph{Proc. NeurIPS}, 2023.

\bibitem{e2tts}
S. E. Eskimez et al., ``E2 TTS: Embarrassingly easy fully non-autoregressive
zero-shot TTS,'' arXiv preprint arXiv:2406.18009, 2024.

\bibitem{f5tts}
Y. Chen et al., ``F5-TTS: A fairytaler that fakes fluent and faithful speech
with flow matching,'' arXiv preprint arXiv:2410.06885, 2024.

\bibitem{valle}
C. Wang et al., ``Neural codec language models are zero-shot text to speech
synthesizers,'' arXiv preprint arXiv:2301.02111, 2023.

\bibitem{hifigan}
J. Kong, J. Kim, and J. Bae, ``HiFi-GAN: Generative adversarial networks for
efficient and high fidelity speech synthesis,'' in \emph{Proc. NeurIPS}, 2020.

\bibitem{fleurs}
A. Conneau et al., ``FLEURS: Few-shot learning evaluation of universal
representations of speech,'' in \emph{Proc. SLT}, 2022.

\bibitem{vctk}
C. Veaux, J. Yamagishi, and K. MacDonald, ``CSTR VCTK corpus: English
multi-speaker corpus for CSTR voice cloning toolkit,'' University of Edinburgh,
2017.

\bibitem{libritts}
H. Zen, V. Dang, R. Clark, Y. Zhang, R. J. Weiss, Y. Jia, Z. Chen, and Y. Wu,
``LibriTTS: A corpus derived from LibriSpeech for text-to-speech,'' in
\emph{Proc. Interspeech}, 2019.

\bibitem{librittsr}
Y. Koizumi et al., ``LibriTTS-R: A restored multi-speaker text-to-speech
corpus,'' in \emph{Proc. Interspeech}, 2023.

\bibitem{whisper}
A. Radford et al., ``Robust speech recognition via large-scale weak
supervision,'' in \emph{Proc. ICML}, 2023.

\bibitem{phowhisper}
V.-A. Nguyen et al., ``PhoWhisper: Automatic speech recognition for
Vietnamese,'' arXiv preprint arXiv:2406.02555, 2024.

\bibitem{utmos}
T. Saeki, D. Xin, W. Nakata, T. Koriyama, S. Takamichi, and H. Saruwatari,
``UTMOS: UTokyo-SaruLab system for VoiceMOS Challenge 2022,'' in
\emph{Proc. Interspeech}, 2022.
\end{thebibliography}

\end{document}
